% ==========================================
% CHAPTER 4
% ==========================================
\chapter{Conditionals and Loops}

\section{Conditional Statements}
As you might have guessed from the chapter title, we create programs like this by implementing various conditional statements and loops.

\textbf{Definition 4.1.1} A conditional statement begins a defined, separated block of code which only executes (runs) if the conditional statement is evaluated by the interpreter to be ``true''.

\textbf{Example 4.1 A Simple Conditional}
\begin{lstlisting}
x = 5
y = 7
if 2*x**2 > y**2:
    print('Wow, thats cool!')
\end{lstlisting}

\subsection{Combining Conditionals}
We are not limited to one conditional per statement; we can combine as many as we need (within reason).

\textbf{Exercise 4.1 Multiple Conditionals}
\begin{lstlisting}
x = input('Enter a number: ')
x = float(x)
y = 15
z = 20

if (x > y) and (x != z):
    print('Nice!')
if (z > x) or (x != y):
    z = x + y + z
\end{lstlisting}

\textbf{Example 4.2 An Else Statement}
\begin{lstlisting}
x = input('Enter a number: ')
if int(x) == 5:
    print('Wow, this was an unlikely coincidence.')
else:
    print('Well, that is interesting.')
\end{lstlisting}

\section{Loops}
The two primary loops in Python are the \texttt{while} and \texttt{for} loops:

\textbf{Definition 4.2.1} A while-loop is a set off block of code that will continue to run sequentially, over and over, so long as a certain condition is met.

\textbf{Definition 4.2.2} A for-loop is a set off block of code that contains a temporary variable known as an iterator, and runs the block of code over and over for different specified values of that iterator.

\subsection{While-Loops}
Lets begin with a simple example of a while-loop.

\textbf{Example 4.3 A while-loop}
\begin{lstlisting}
x = 100
while x > 5:
    print(x)
    x = x - 1
\end{lstlisting}

\subsection{For-Loops}
For-loops are one of the most powerful tools in Python. What they allow us to do is write a block of code that's like a template- it has the code we want to run, but without defining exactly ``on what'' the code acts.

\textbf{Exercise 4.3 A simple for-loop}
\begin{lstlisting}
arr = [1,2,3,4,5,6,7,8,9,10]
for i in arr:
    if i % 2 == 0:
        print(i)
\end{lstlisting}

\subsection{Nested For-Loops}
Just briefly, we'd like to mention that you can in fact nest multiple for-loops together, if you need to iterate over more than one value in your code. This often happens when dealing with two-dimensional arrays.

\textbf{Example 4.4 Iterating a 2D Array}
\begin{lstlisting}
for i in range(len(x)-1):
    for j in range(len(y)-1):
        if arr[i,j] < 1500.0:
            arr[i,j] = 0
\end{lstlisting}